\documentclass{article}
\usepackage[utf8]{inputenc}
\usepackage[fa-IR]{polyglossia}
\setmainlanguage{persian}
\setotherlanguage{english}
\newfontfamily\persianfont[Script=Arabic]{Arial} % Or any Persian font available in Overleaf

\title{گزارش پروژه پایش بحران WarIran - فاز ۱}
\author{Antigravity AI Assistant}
\date{\today}

\begin{document}
\maketitle

\section{مقدمه}
پروژه WarIran با هدف ایجاد یک سیستم هوشمند برای پایش مناطق مختلف ایران با استفاده از داده‌های ماهواره‌ای و حرارتی کلید خورد. در این فاز، زیرساخت اصلی پروژه و اولین موتور تشخیص (حرارتی) پیاده‌سازی شده است.

\section{اقدامات انجام شده}
\begin{itemize}
    \item \textbf{زیرساخت برنامه‌نویسی:} راه‌اندازی ساختار ماژولار پروژه با استفاده از زبان پایتون.
    \item \textbf{استانداردسازی (Docker):} استفاده از Docker برای تضمین اجرای یکسان برنامه در تمامی محیط‌ها.
    \item \textbf{امنیت داده‌ها:} پیاده‌سازی سیستم مدیریت کلیدهای API با استفاده از فایل .env.
    \item \textbf{موتور تشخیص حرارتی:} اتصال به سرویس NASA FIRMS برای شناسایی نقاط حرارتی (انفجار/آتش‌سوزی) در زمان واقعی.
\end{itemize}

\section{وضعیت فعلی}
در حال حاضر موتور حرارتی قادر است با دقت بالا، مناطق با مختصات مشخص را از نظر تغییرات دمایی ناگهانی بررسی کرده و در صورت بروز حادثه، داده‌های مربوط به موقعیت جغرافیایی را استخراج کند.

\section{گام‌های بعدی}
\begin{enumerate}
    \item اتصال به Google Earth Engine برای دریافت تصاویر راداری Sentinel-1.
    \item پیاده‌سازی الگوریتم تشخیص خرابی‌های فیزیکی (Structural Damage).
    \item راه‌اندازی رابط کاربری وب و سیستم اعلان تلگرام.
\end{enumerate}

\end{document}
